\section{کنترل‌گر مدیریت برد (\lr{BMC}) و راه‌اندازی}

\subsection*{بخش آ – مسیر راه‌اندازی میزبان و جایگاه \lr{BMC}}

برنامهٔ راه‌اندازی (میان‌افزار، در \lr{ROM/EPROM}) هنگام روشن شدن یا راه‌اندازی دوباره اجرا می‌شود، همه‌چیز را مقداردهی می‌کند (ثبات‌های \lr{CPU}، کنترل‌گرهای دستگاه)، هسته را بارگذاری می‌کند و هسته نخستین پردازه را اجرا می‌کند؛ مثال لینوکس: \lr{BIOS $\rightarrow$ MBR bootloader $\rightarrow$ GRUB $\rightarrow$ Kernel init $\rightarrow$ INIT $\rightarrow$ Run level}. مراحل \lr{UEFI} هم \lr{SEC}، \lr{PEI}، \lr{DXE} و \lr{BDS} است. \lr{BMC} رایانهٔ کوچکی (معمولاً \lr{SoC} مبتنی بر \lr{ARM}) روی برد اصلی است که \textbf{در مسیری جدا} قرار دارد.

\begin{center}
\begin{tikzpicture}[
  font=\small, >=Stealth,
  box/.style={draw=navy, fill=lane1, rounded corners=3pt, minimum height=1.05cm, minimum width=5.7cm, align=center, inner sep=3pt, line width=0.6pt},
  bmc/.style={draw=accent, fill=lane2, rounded corners=3pt, minimum height=1.05cm, minimum width=5.7cm, align=center, inner sep=3pt, line width=0.6pt},
  fin/.style={draw=navy, fill=white, rounded corners=3pt, minimum height=1.05cm, minimum width=5.7cm, align=center, inner sep=3pt, line width=1.2pt, dashed},
  main/.style={->, line width=0.7pt, color=navy},
  side/.style={->, line width=0.7pt, color=accent, dashed},
  lbl/.style={font=\scriptsize, align=center, fill=white, inner sep=1.5pt}
]
\node[box] (h0) at (4.4, 0)    {\RL{\textbf{روشن شدن برق}}\\[-1pt]\RL{\scriptsize بررسی سلامت برق (\lr{power-good})}};
\node[box] (h1) at (4.4,-1.65) {\RL{\textbf{میان‌افزار \lr{UEFI/BIOS}}}\\[-1pt]\RL{\scriptsize \lr{SEC}، \lr{PEI}: مقداردهی \lr{CPU}، حافظه و تراشه‌ها}};
\node[box] (h2) at (4.4,-3.30) {\RL{\textbf{\lr{DXE} و \lr{BDS}}}\\[-1pt]\RL{\scriptsize بارگذاری درایورها، انتخاب دستگاه بوت}};
\node[box] (h3) at (4.4,-4.95) {\RL{\textbf{بوت‌لودر (\lr{MBR/GRUB})}}\\[-1pt]\RL{\scriptsize پیدا کردن و بارگذاری هسته در حافظه}};
\node[box] (h4) at (4.4,-6.60) {\RL{\textbf{هسته (\lr{Kernel init})}}\\[-1pt]\RL{\scriptsize راه‌اندازی زیرسیستم‌ها و سوار کردن ریشه}};
\node[box] (h5) at (4.4,-8.25) {\RL{\textbf{نخستین پردازه (\lr{init/PID 1})}}\\[-1pt]\RL{\scriptsize اجرای سرویس‌ها و پردازه‌های کاربر}};
\node[fin] (h6) at (4.4,-9.90) {\RL{\textbf{محیط اجرای پردازه‌های کاربر}}\\[-1pt]\RL{\scriptsize (هدف مسئله: میزبان به اینجا نمی‌رسد)}};
\foreach \a/\b in {h0/h1,h1/h2,h2/h3,h3/h4,h4/h5,h5/h6} \draw[main] (\a.south) -- (\b.north);
\node[bmc] (b0) at (-4.4, 0)    {\RL{\textbf{برق آماده‌به‌کار و بوت داخلی}}\\[-1pt]\RL{\scriptsize \lr{BMC} پیش از میزبان و مستقل از آن بالا می‌آید}};
\node[bmc] (b1) at (-4.4,-1.65) {\RL{\textbf{رابط مدیریتی}}\\[-1pt]\RL{\scriptsize \lr{IPMI/Redfish/Web} روی شبکهٔ اختصاصی یا مشترک}};
\node[bmc] (b2) at (-4.4,-3.30) {\RL{\textbf{نمایش وضعیت و ثبت فرمان‌ها}}\\[-1pt]\RL{\scriptsize وضعیت میزبان، بازنشانی، حسگرها، کنسول}};
\foreach \a/\b in {b0/b1,b1/b2} \draw[main,color=accent] (\a.south) -- (\b.north);
\draw[side] (b2.east) -- ++(1.55,0) |- (h0.west);
\node[lbl] at (0,-1.65) {\RL{فرمان روشن،}\\\RL{خاموش، بازنشانی}\\\RL{(خط سخت‌افزاری)}};
\draw[side] (h4.west) -| (b2.south);
\node[lbl] at (-1.4,-6.60) {\RL{وضعیت، حسگرها،}\\\RL{کنسول سریال}};
\node[font=\bfseries\small, color=navy] at (4.4, 0.95) {\RL{مسیر اصلی (میزبان)}};
\node[font=\bfseries\small, color=accent] at (-4.4, 0.95) {\RL{مسیر مدیریتی جداگانه (\lr{BMC})}};
\end{tikzpicture}
\end{center}

\textbf{توضیح تصویر:} مسیر اصلی (راست) از برق، میان‌افزار، بوت‌لودر و هسته به نخستین پردازه می‌رسد. \lr{BMC} (چپ) رایانهٔ مستقل با برق آماده‌به‌کار، میان‌افزار و رابط شبکهٔ خودش است و فقط از طریق خطوط سخت‌افزاری و پیوندهای جانبی (روشن/بازنشانی، حسگرها، کنسول سریال) با مسیر میزبان ارتباط دارد. به همین دلیل در تمام مراحل و حتی وقتی میزبان به محیط اجرای پردازه‌های کاربر نمی‌رسد، فعال می‌ماند.

\subsection*{بخش ب – آیا رابط مدیریتی \lr{BMC} پردازه‌ای در سیستم‌عامل میزبان است؟}

\textbf{خیر.} دلیل‌ها:
\begin{itemize}
  \item \lr{BMC} یک رایانهٔ کوچک مستقل روی برد اصلی است (معمولاً \lr{ARM-based SoC}) با \lr{CPU}، حافظه و میان‌افزار (و معمولاً سیستم‌عامل) خودش. پردازه در مفهوم درس یعنی برنامه‌ای که هستهٔ میزبان آن را ایجاد و زمان‌بندی می‌کند و از حافظهٔ میزبان استفاده می‌کند؛ رابط \lr{BMC} هیچ‌کدام را ندارد.
  \item برق آماده‌به‌کار و درگاه شبکهٔ اختصاصی (یا مشترک از طریق \lr{NC-SI}) دارد؛ پس حتی وقتی میزبان خاموش است یا هنوز هسته‌ای بارگذاری نشده، کار می‌کند. چیزی که پیش از وجود سیستم‌عامل کار می‌کند نمی‌تواند پردازهٔ همان سیستم‌عامل باشد.
  \item بر اساس صورت مسئله، رابط مدیریتی با مسیر مدیریتی جداگانه از میزبان است و حتی وقتی سیستم‌عامل میزبان اجرا نشده پاسخ می‌دهد. بنابراین نه از هستهٔ میزبان زمان می‌گیرد و نه به آن وابسته است.
\end{itemize}

\subsection*{بخش ج – آیا پاسخ‌گو بودن \lr{BMC} در حالی که میزبان به محیط پردازه‌ها نرسیده تناقض ایجاد می‌کند؟}

\textbf{خیر، تناقضی نیست}، چون \lr{BMC} و میزبان دو سامانهٔ جدا هستند.
\begin{itemize}
  \item \textbf{آنچه می‌توان نتیجه گرفت:} برق آماده‌به‌کار برقرار است، \lr{BMC} بوت شده و مسیر مدیریتی (شبکه و رابط) سالم است، و \lr{BMC} می‌تواند وضعیت سخت‌افزاری برد را بخواند (برق، حسگرها) و فرمان بازنشانی را ثبت کند.
  \item \textbf{آنچه نمی‌توان نتیجه گرفت:} \lr{BMC} اطلاعات خود را از سیگنال‌ها و حسگرهای سخت‌افزاری می‌گیرد؛ مثلاً می‌بیند میزبان «روشن است». این به معنای سلامت یا اجرای میان‌افزار، بوت‌لودر، هسته یا هر پردازه‌ای نیست. نمی‌توان نتیجه گرفت که میزبان به‌درستی بوت شده یا از کار نیفتاده است؛ فقط می‌توان گفت میزبان به محیط اجرای پردازه‌های کاربر نرسیده است و علت آن (میان‌افزار، بوت‌لودر، هسته یا سخت‌افزار) را باید جای دیگر پیدا کرد.
\end{itemize}

\subsection*{بخش د – بارگذاری هسته پس از ثبت درخواست بازنشانی}

پس از روشن شدن دوباره، ابتدا میان‌افزار (\lr{UEFI/BIOS}) اجرا می‌شود و سخت‌افزار را مقداردهی می‌کند، سپس بوت‌لودر (\lr{GRUB}) \textbf{هسته} را از دیسک در حافظه بارگذاری می‌کند و کنترل را به آن می‌دهد. \textbf{ثبت درخواست بازنشانی به معنای اجرای سیستم‌عامل میزبان توسط \lr{BMC} نیست}، زیرا:
\begin{itemize}
  \item \lr{BMC} فقط یک «فرمان» ثبت و ارسال می‌کند (فعال‌کردن خط بازنشانی یا چرخهٔ برق؛ معادل زدن دکمهٔ \lr{reset}) و کد یا حافظهٔ میزبان را اجرا نمی‌کند.
  \item پس از دریافت سیگنال، مسئولیت اجرا با سخت‌افزار و میان‌افزار میزبان است که از ابتدا شروع می‌کنند و فرایند بوت را طی می‌کنند. کد \lr{BMC} روی پردازندهٔ جداگانه‌ای اجرا می‌شود و اصلاً به میزبان وارد نمی‌شود.
  \item هسته‌ای که بارگذاری می‌شود همان است که روی دیسک میزبان قرار دارد و \lr{BMC} در انتخاب یا بارگذاری آن نقشی ندارد (مگر فرمان‌هایی مثل تعیین دستگاه بوت در دفعهٔ بعد، که باز هم فقط یک تنظیم برای میان‌افزار است).
\end{itemize}

\subsection*{بخش ه – میزبان به درخواست‌های پردازهٔ کاربر پاسخ می‌دهد، ولی رابط \lr{BMC} در دسترس نیست}

\textbf{نمی‌توان نتیجه گرفت که سیستم‌عامل میزبان از کار افتاده است.} سیستم‌عامل میزبان به \lr{BMC} وابسته نیست؛ پاسخ‌گو بودن پردازه‌های کاربر نشان می‌دهد که هسته، زمان‌بندی و مسیر شبکهٔ میزبان سالم است. نبودن رابط \lr{BMC} به یکی از این علت‌ها می‌تواند باشد: \lr{BMC} در حال راه‌اندازی دوباره یا به‌روزرسانی میان‌افزار است، خود \lr{BMC} (میان‌افزار یا سخت‌افزارش) خراب شده، شبکهٔ مدیریتی یا درگاه اختصاصی (یا مسیر \lr{NC-SI}) قطع است، یا تنظیمات دسترسی عوض شده است. یعنی خرابی \lr{BMC} فقط توانایی مدیریت از راه دور را کم می‌کند و نشانهٔ خرابی میزبان نیست.

\subsection*{بخش و – مقایسهٔ دو حالت}

\begin{center}\small
\small
\begin{tabular}{@{}>{\raggedright\arraybackslash}p{4.4cm}>{\raggedright\arraybackslash}p{5.3cm}>{\raggedright\arraybackslash}p{5.3cm}@{}}
\toprule
  {\textbf{حالت}} &
  {\textbf{مسیری که کار می‌کند}} &
  {\textbf{مسیری که ممکن است از کار افتاده باشد}} \\
\midrule
  {\lr{BMC} در دسترس، سیستم‌عامل میزبان در دسترس نیست} &
  {مسیر برون‌باند \lr{BMC}: دیدن وضعیت برق و حسگرها، روشن/خاموش‌کردن و بازنشانی، کنسول سریال از راه دور، ثبت رویدادها} &
  {مسیر درون‌باند میزبان: هسته، پردازه‌های کاربر، سرویس‌هایی مثل \lr{SSH}، شبکهٔ میزبان؛ (سیستم‌عامل پایین است یا اصلاً بالا نیامده)} \\
\addlinespace
  {\lr{BMC} در دسترس نیست، سیستم‌عامل میزبان در دسترس است} &
  {مسیر درون‌باند میزبان: سیستم‌عامل، پردازه‌ها و شبکهٔ میزبان؛ (به‌جز حالتی که کارت شبکه را مشترک با \lr{BMC} استفاده کنند و آن مسیر هم قطع شود)} &
  {مسیر مدیریتی برون‌باند: رابط \lr{BMC}، کنترل از راه دور برق و بازنشانی، کنسول، حسگرها؛ (برق و میزبان مشکلی ندارند)} \\
\bottomrule
\end{tabular}
\end{center}
